\documentclass[11pt,letterpaper]{article}
\usepackage[top=.88in,bottom=.86in,left=.94in,right=.94in,headheight=14pt,headsep=16pt]{geometry}
\usepackage[T1]{fontenc}\usepackage{lmodern}
\usepackage{amsmath,amssymb,amsthm,booktabs,array,xcolor,fancyhdr,titlesec,hyperref}
\definecolor{ink}{HTML}{19384C}\definecolor{teal}{HTML}{08736E}
\hypersetup{colorlinks=true,linkcolor=teal,citecolor=teal,urlcolor=teal,
 pdftitle={Beyond Dirichlet GRH: four transfers from a Lean proof},pdfauthor={Research report prepared for Vladimir Reshetnikov}}
\pagestyle{fancy}\fancyhf{}\lhead{\small\sffamily\color{ink}Beyond Dirichlet GRH}
\rhead{\small\sffamily\color{ink}Mathematical research report}\cfoot{\small\thepage}
\renewcommand{\headrulewidth}{.3pt}
\titleformat{\section}{\Large\sffamily\bfseries\color{ink}}{\thesection}{.65em}{}
\titleformat{\subsection}{\large\sffamily\bfseries\color{ink}}{\thesubsection}{.6em}{}
\newtheorem{theorem}{Theorem}[section]
\newcommand{\source}[2]{\href{https://github.com/lsernau/grh-lean/blob/72fead102f663d5c7a407964a13dc0cf9d977207/grh/src/#1}{#2}}
\newcommand{\prev}[3]{\href{https://github.com/VladimirReshetnikov/ProveIt/blob/#1/docs/incoming/#2}{#3}}
\setlength{\emergencystretch}{2em}\setlength{\parskip}{2pt}
\raggedbottom
\begin{document}
\thispagestyle{plain}
\noindent{\sffamily\small\color{teal}\bfseries RESEARCH REPORT \hfill 26 NOVEMBER 2026}\par
\vspace{6pt}
\noindent{\Huge\sffamily\bfseries\color{ink}Beyond Dirichlet GRH}\par
\vspace{5pt}
\noindent{\large Four applications of progression avoidance, cyclotomic valuations,\\
half-plane exclusion and sharpness at finite precision}\par
\vspace{10pt}
\noindent The Lean proof contains a finite construction that avoids many weighted means simultaneously and precise calculations at roots of unity. These ingredients can be used independently of its theorem on $L$-function zeros. The finite-ring and graph-tower results answer further questions in the ProveIt reports. Exact positivity certificates locate the zeros of infinitely many Ehrhart polynomials. A local formal-group calculation determines when an earlier interpolation bound is attained.

\medskip
\begin{center}\small
\begin{tabular}{@{}>{\raggedright\arraybackslash}p{.25\linewidth}>{\raggedright\arraybackslash}p{.68\linewidth}@{}}
\toprule
\textbf{Problem} & \textbf{Result proved here}\\ \midrule
Geometric progressions over finite rings (\S\ref{sec:rings}) &
An exact convergence criterion for growing products of local rings with square-zero maximal ideal. In odd characteristic the single-factor loss also has an exact leading constant.\\[4pt]
Spanning trees in graph towers (\S\ref{sec:rank-two-onset}) &
An explicit logarithmic bound in the voltage height for rank two. For fixed prime $p$ and loop count $m$, worst onset grows linearly in the bit length exactly when $m\ge\max(3,p)$.\\[4pt]
Ehrhart zeros (\S\ref{sec:ehrhart}) &
For every complete bipartite graph whose smaller part has size at most 33, all roots are simple and lie on the canonical line. Through width 32 we certify interlacing in both part-size directions.\\[4pt]
Lubin--Tate interpolation (\S\ref{sec:branch-equality}) &
A full equality criterion for the pole-order bound. An exact pole-order formula holds at any precision for contact height one, and at precision two for all contact heights.\\
\bottomrule
\end{tabular}
\end{center}
\medskip
\noindent\textbf{Density loss at fixed residue field.} Let $A$ have fixed residue field of order $q$, cardinality $Q$ and square-zero maximal ideal. For ordinary three-term progressions let $\Delta(A)$ denote the supremum of upper densities of admissible sets. As $Q\to\infty$, we obtain
\[
  1-\Delta(A)=
  \begin{cases}
      Q^{-2}+O(Q^{-3}),&q\text{ even},\\
      (1-q^{-q})/(qQ)+O_q(Q^{-2}),&q\text{ odd}.
  \end{cases}
\]
The odd constant comes from residue polynomials divisible by every monic linear polynomial. They have proportion $q^{-q}$ and require a different treatment in the packing argument (p.~\pageref{sec:odd-leading}).

\medskip
\noindent\textbf{Status.} The arguments below are mathematical proofs with executable finite checks. The positivity claims in the Ehrhart proof use exact integer calculations as part of the proof. The new results have not been formalized in Lean or externally refereed. Priority has not been established. The role of each earlier report or paper is identified where the corresponding result is stated.
\clearpage
\section{What can be taken out of the Lean proof?}
At commit \texttt{72fead102f663d5c7a407964a13dc0cf9d977207}, the readable project proves Mathlib's Riemann hypothesis and four related statements. The right-half-plane nonvanishing statement quantifies over every Dirichlet character of positive modulus, with the principal-character pole excepted. The critical-strip statement includes imprimitive characters. The exported scope is Dirichlet $L$-functions. The project pins Lean 4.35.0 and Mathlib \texttt{56a00585}. Its authors report kernel and independent export checks with the usual three axioms. For this report we inspected the source at that revision. The upstream proof was not rebuilt here. \source{Statements.lean}{Statements}, \source{Main.lean}{public theorems}, \cite{lean}.

\paragraph{The arithmetic step.}
The proof compares two estimates for the same family of target forms. Short integral forms give an upper bound for a Gram determinant and a local index. A second-moment argument supplies sufficiently many independent images. A lower bound for the integral forms, using twisted determinants and height estimates, contains the same large main term. Its cancellation leaves a one-sided bound on an arithmetic remainder. A choice of quadratic twist removes twisted zeros in the relevant box. The residue of a transform then records a sum of nonnegative orders of untwisted zeros. A one-sided Landau argument forces that sum to vanish in the half-plane being tested. The functional equation gives the critical-line conclusion. Extending this argument to another analytic family would require a suitable remainder estimate and a way to isolate the desired poles. These are substantive arithmetic inputs. \source{New/Reduction/MainInequality.lean}{Determinant comparison}, \source{New/Reduction/Zeros/Criterion.lean}{one-sided criterion}, \source{New/Reduction/OneSidedBound.lean}{arithmetic bound}.

\paragraph{How each application uses the source.}
The finite-ring lower bound applies the weighted Behrend lemma from the Lean source, with the weights now representing unequal gaps between exponents of a polynomial ratio. The graph theorem uses the same cyclotomic specialization identity and adds a separation argument for projected exponents. For Ehrhart polynomials the transfer is methodological: rotation and an exclusion in one half-plane suffice by symmetry, but the needed positivity comes from the known Ehrhart expansion and new finite certificates. Finally the interpolation result pursues the local sharpness question identified in an earlier report inspired by the source's derivative-extremizer construction. Its proof uses formal groups and Laurent-series arithmetic. Each application stands as a separate argument. \source{New/Reduction/Data/Existence.lean}{Weighted means}, \source{New/Congruences/Selmer/RootsOfUnity.lean}{torsion valuations}, \source{New/Derivatives/Sharpness/OneCurve.lean}{derivative extremizers}.

\paragraph{Earlier results used for comparison.}
The 55 ZIPs in the incoming directory at revision \texttt{bd8864aa31} were searched for preceding results. Its polynomial-density reports already give improvements over the greedy constructions and a counterexample to the proposed upper-density value over fields. Our finite-ring application starts from their question about nilpotents. In the graph problem, eventual polynomiality and even a terminating determination of the first valid level were available; the question is how late stabilization can occur as the input grows. The interpolation report already identified the individual branch cutoff as one plus a pole order. Our contribution there is to compute its loss through cancellation. The linked packets and primary sources give the precise comparisons. \cite{products,graphprior,branchprior,fields}

\clearpage
\section{Progression-free sets over finite local rings}\label{sec:rings}

Asada, Fourakis, Manski, McNew, Miller and Moreland ask how geometric-progression avoidance over $\mathbb F_q[x]$ changes when another finite ring replaces the coefficient field \cite[Question 6.5]{asada}. An earlier ProveIt manuscript classifies growing products of fields using disjoint triples and the density Hales--Jewett theorem. It ends by asking for a classification with nonzero radicals \cite{products}. We resolve the growing-product question for all finite commutative local rings whose maximal ideals square to zero, including mixed characteristic. A bounded family of unequal exponent gaps has exactly the same threshold.

\subsection{Statement and consequences}
All rings have an identity, and all local rings here have a nonzero residue field. Put $B_n(R)=\{f\in R[x]:\deg f\le n\}$, including zero, and take $\deg 0=-\infty$. For an integer $L\ge2$, call a set admissible if it contains no three distinct polynomials of the form
\begin{equation}\label{eq:pattern}
              f,\quad r^a f,\quad r^{a+b}f,\qquad
 r\in R[x]\text{ a nonunit},\quad a,b\ge1,\quad a+b\le L .
\end{equation}
Thus $L=2$ is the published progression problem. Write $\Delta_L(R)$ for the supremum over admissible $E$ of $\limsup_{n\to\infty}|E\cap B_n(R)|/|R|^{n+1}$. The proofs work also if triples with repeated terms are forbidden. In lower-bound constructions we discard polynomials with a zero residue component, and all triples used in upper bounds are distinct.

\begin{theorem}\label{thm:local}
Let $(A_i,\mathfrak m_i)$ be finite commutative local rings with $\mathfrak m_i^2=0$. Write $q_i=|A_i/\mathfrak m_i|$, $Q_i=|A_i|$, and $R_s=\prod_{i=1}^s A_i$. Define
\[
 w_i=\begin{cases}
       Q_i^{-2},& q_i\text{ even},\\
       (q_iQ_i)^{-1},& q_i\text{ odd}.
     \end{cases}
\]
For every fixed $L\ge2$ the limit $\lim_{s\to\infty}\Delta_L(R_s)$ exists, and it is positive if and only if $\sum_i w_i<\infty$. For one local factor, uniformly in $L$,
\[
       \gamma_i w_i\ \le\ 1-\Delta_L(A_i)\ \le\ w_i,\qquad
 \gamma_i=\begin{cases}
    1-Q_i^{-1},&q_i\text{ even},\\
    (1-q_i^{-1})^2,&q_i\text{ odd}.
 \end{cases}
\]
\end{theorem}

Here $Q_i$ and $q_i$ measure different features. At fixed odd $q_i=q$ the density loss has order $Q_i^{-1}$, whereas at fixed even $q$ it has order $Q_i^{-2}$. For example take $A_i=\mathbb F_q\oplus V_i$ with $V_i^2=0$ and $1+\dim V_i=\max(2,\lceil\log_q(i+1)\rceil)$. Then $Q_i$ is comparable to $i+1$ for fixed $q$. The limiting extremal density is positive for $q=2$ and zero for $q=3$. Both families have growing square-zero radicals. Also, $A=\mathbb Z/p^2\mathbb Z$ is covered with $q=p$ and $Q=p^2$.

\subsection{The degree restrictions at a single factor}
Write $\bar f$ for coefficientwise reduction and $k=A/\mathfrak m$. A polynomial is a unit if and only if its reduction is a nonzero constant: lift that constant to a unit and invert the remaining nilpotent correction. The same fact holds for every finite local ring since its maximal ideal is nilpotent. In the present square-zero case, a polynomial with reduction of degree $e\ge1$ can be written $r=t+u$ with $\deg t=e$, leading coefficient of $t$ a unit, and $u\in\mathfrak m[x]$. Choose lifts only of the coefficients of $\bar r$ up to degree $e$. Then
\[
              r^j=t^j+j t^{j-1}u,\qquad j\ge1 .
\]
If the residue characteristic $p$ divides $j$, then $ju=0$, since $p\in\mathfrak m$. Multiplication by $t$ increases the degree of every nonzero polynomial by exactly $e$. These assertions allow the characteristic of $A$ itself to be $p^2$.

Suppose $f,y=r^af,z=r^{a+b}f$ are in $B_n(A)$ with all three reductions nonzero. If $r$ is a nonunit, its reduction must have positive degree $e$. In residue characteristic two, at least one of $a,b,a+b$ is even. The corresponding relation $y=t^a f$, $z=t^b y$, or $z=t^{a+b}f$ forces $f$ or $y$ to have degree at most $n-2$. Remove $B_{n-2}(A)$, of relative size $Q^{-2}$.

For odd $p$ remove instead
\[
        D_n=\{f:\deg f\le n-1,\ \deg\bar f\le n-2\}.
\]
Its probability is $(qQ)^{-1}$ by the independent uniform coefficients of $x^n$ and $x^{n-1}$. If $p\mid b$, then $z=t^b y$ forces $\deg y\le n-2$. Otherwise $b$ is a unit in $A$, and elimination of $u$ gives
\begin{equation}\label{eq:eliminate}
       (a+b)t^b y-a z=b t^{a+b} f .
\end{equation}
Consequently $\deg f\le n-ae\le n-1$. Also $\bar z=\bar r^{a+b}\bar f$ forces $\deg\bar f\le n-(a+b)e\le n-2$. The first term lies in $D_n$. This proves the asserted finite-ball deletions. Removing the zero-reduction fiber as well costs at most $q^{-(n+1)}$.

\subsection{Disjoint triples and the odd-characteristic construction}
One can pack triples occupying a positive multiple of the deleted fraction. By a packing we mean a collection of vertex-disjoint ordered triples \eqref{eq:pattern} with $a=b=1$. It is enough to specify their first terms and an injective rule for each of the three positions.

For even $q$ take every $f$ of actual degree $n-2$, with triple $(f,xf,x^2f)$. The three positions have different actual degrees. The proportion of first terms is $(1-Q^{-1})Q^{-2}$. This argument also applies over arbitrary finite coefficient rings.

For odd $q$ and $n\ge3$ take all $f=\sum f_jx^j$ satisfying
\[
   f_n=0,\qquad f_{n-1}\in\mathfrak m,\qquad
   f_{n-2},\,f_0\in A^\times .
\]
There are $(1-q^{-1})^2/(qQ)$ of these as a proportion of the whole ball. Set
\[
       c=-\frac{f_{n-1}}{2 f_{n-2}}\in\mathfrak m,\qquad
       r=x+c(x^2+1),\qquad y=rf,\quad z=r^2f .
\]
The term $x^2c$ cancels the unwanted degree $n+1$ term of $z$. Indeed $c f_{n-1}=0$ and $f_{n-1}+2c f_{n-2}=0$; both $y$ and $z$ are in $B_n(A)$. Their residue degrees are $n-1,n$, distinct from the starting degree $n-2$ in the residue field. To recover the start, the scalar polynomial $\bar f$ is obtained from $\bar y=x\bar f$ or $\bar z=x^2\bar f$. The coefficients
\[
       [x^n]y=-c f_{n-2},\qquad [x]z=2c f_0
\]
determine $c$ from either position: $\mathfrak m$ is a vector space over $k$ and the indicated scalars are known and nonzero in $k$. The constant correction $c$ in $r$ is what permits recovery at the third position. Finally $r(x-c(x^2+1))=x^2$ shows that multiplication by $r$ and by $r^2$ is injective. The triples are pairwise disjoint, including between different starting polynomials. This proves the single-factor upper estimate on density and supplies the product obstruction.

\subsection{From finite balls to arbitrarily many factors}
For a fixed finite product let $M(n)$ be the largest fraction of its ball occupied by an admissible set. We record why $\Delta_L(R_s)=\limsup_n M(n)$. Choose a subsequence $n_j$ approaching the upper limit so rapidly that
\[
      n_j-j>(1+L)n_{j-1},\qquad n_j-(1+L)j>n_{j-1}.
\]
From an optimizer at $n_j$ keep only elements whose residue degrees in all coordinates lie in $[n_j-j,n_j]$. The fraction discarded is at most $\sum_{i\le s}q_i^{-j-1}$, by uniform reduction. The union of the retained blocks is admissible. In each coordinate the three residue degrees are $d,d+ae,d+(a+b)e$ for $e\ge0$, and all are finite. If just the last term were in the latest block, its degree in each coordinate would be at most $(1+L)n_{j-1}$ because the middle degree is at most $n_{j-1}$. If the last two were in that block, their difference $be$ would be at most $j$, so the first degree would be at least $n_j-(1+L)j$. Both possibilities contradict the separation. This proves the equality with $\limsup M(n)$; the converse inequality is immediate. Slicing the product by its last coordinate shows $\Delta_L(R_{s+1})\le\Delta_L(R_s)$, using ratio $1$ in the coordinate held fixed.

The product of the single-factor constructions, trimmed to nonzero residues, gives
\begin{equation}\label{eq:product-bound}
                    \Delta_L(R_s)\ge \prod_{i=1}^s(1-w_i).
\end{equation}
In a progression the ratio must be a nonunit in at least one coordinate. There all three residue degrees are finite, so the deletion argument in that coordinate applies.

For the reverse bound put $u_i=\gamma_i w_i$. Partition $B_n(A_i)$ into the packed triples above and singleton cells. A uniformly drawn point falls in a triple with probability $3u_i$, independently across the coordinates. Conditional on the cells, if $K$ coordinates fall in triples the point is uniform on a ternary cube of dimension $K$. Every combinatorial line gives a forbidden triple, since its moving coordinates have their respective fixed ratios, and the other coordinates can use ratio $1$. If $\alpha_k$ denotes the maximal density of a line-free subset of $\{0,1,2\}^k$, with $\alpha_0=1$, then
\[
          M(n)\le \mathbb E\,\alpha_K,\qquad
          K=\sum_{i=1}^s X_i,\quad X_i\sim{\rm Bernoulli}(3u_i)
          \ \text{independently}.
\]
The density Hales--Jewett theorem gives $\alpha_k\to0$ \cite{dhj}. Since $\gamma_i\ge4/9$ and $\operatorname{Var}K\le\mathbb EK$, divergence of $\sum_iw_i$ gives $K\to\infty$ in probability. The upper bound then tends to zero. When the sum converges, the product in \eqref{eq:product-bound} is positive, since $w_i\le1/4$. This completes the proof of Theorem~\ref{thm:local}.

\subsection{A direct use of the weighted Behrend lemma}
There is also a quantitative lower bound valid for \emph{every} finite product of commutative local rings, with radicals of arbitrary nilpotence order. Put $V=\sum_i q_i/(q_i-1)^2$. For a constant $C_L$ depending only on the allowed exponent sum,
\begin{equation}\label{eq:behrend}
            \Delta_L(R_s)\ \ge\ \exp[-C_L\sqrt{\log(2+V)}] .
\end{equation}
The incoming field-product argument treats ordinary midpoints \cite{products}. Here the choices $a,b$ of \eqref{eq:pattern} give precisely the several weights that the Lean construction handles simultaneously.

For each integer $m\ge2$, the lemma \texttt{exists\_three\_term\_free} gives a set $S\subset[0,m)$ of size at least $m\exp[-C'_L\sqrt{\log(2+m)}]$ avoiding all nonconstant weighted means $(a+b)d_2=b d_1+a d_3$ with $a+b\le L$. Indeed choose dimension $\lfloor\sqrt{\log m}\rfloor$ and digit-side length $\lfloor m^{1/d}/L\rfloor$ for sufficiently large $m$; absorb rounding and smaller $m$ in the constant. The Lean lemma obtains a whole congruence-free set, stronger than the integer version needed here \cite{lean}.

At cutoff $n$ sum the residue degree deficits $D_n=\sum_i(n-\deg\bar f_i)$ over tuples with each residue nonzero. Its limiting law for fixed $s$ is the sum $D$ of independent variables with masses $(1-q_i^{-1})q_i^{-j}$ for $j\ge0$. Its variance is $V$. By Chebyshev an interval of $m\le4\sqrt V+4$ consecutive integers, enlarged to length two if needed, carries at least $3/4$ of the mass. Among the $2m-1$ translations of $S$ that could meet that interval some translate receives probability at least $3|S|/(4(2m-1))$. Use that translate to select $D_n$. Along any surviving progression the sum of the residue degrees increases with step size $\sum_i\deg\bar r_i>0$ before multiplication by the respective exponents $a,b$. Its three deficits would satisfy the forbidden weighted mean. Passing to the limit and applying the block argument proves \eqref{eq:behrend}; it only used finite local rings and reduction to their residue fields.

\subsection{The exact leading constant in odd residue characteristic}
\label{sec:odd-leading}

Write $A$ for a finite commutative local ring with residue field
$k=\mathbb F_q$ of odd characteristic, maximal ideal $J$ satisfying
$J^2=0$, cardinality $Q$, and $M=|J|=Q/q$. Multiplication makes
$J$ a vector space over $k$. Choose a representative $\widetilde\gamma$
in $A$ for each $\gamma\in k$, and put
$P_\gamma=x-\widetilde\gamma$ and
$\Lambda=\prod_{\gamma\in k}(x-\gamma)$ in $k[x]$.
Write $\overline f$ for coefficientwise reduction. Use the same bound $L\ge2$ on the sum of the two exponent gaps as in Theorem~\ref{thm:local}.
Then for each fixed odd $q$,
\[
   \lim_{Q\to\infty} qQ\bigl(1-\Delta_L(A)\bigr)=1-q^{-q}.
\]
The limit holds along any sequence of these rings with fixed residue
field size. In particular the leading constant already accounts for
the positive fraction $q^{-q}$ of residue polynomials divisible by
every monic linear polynomial.

\paragraph{Improved construction.}
At cutoff $n\ge q+3$, consider the rectangle $f_n=0$, $f_{n-1}\in J$;
subscripts indicate coefficients. Its fraction is $1/(qQ)$.
Inside this rectangle delete only the points with
$\Lambda\nmid\overline f$. Their fraction in the full ball is
$(1-q^{-q})/(qQ)$ by uniformity modulo the degree-$q$ polynomial $\Lambda$.
Also delete all points of degree at most $n-2$, and all points divisible
in $A[x]$ by some $P_\gamma^2$. The additional cost is at most
$(q+1)Q^{-2}$ since $P_\gamma$ is monic.
As in the scalar-degree gluing argument, we may remove the
zero-residue points at a cost tending to zero for this fixed ring.

Consider a remaining progression $f,y=r^i f,z=r^{i+j}f$ with a nonunit ratio.
Nonzero reductions at all three terms ensure that
$e_0=\deg\overline r$ is positive.
Write $t$ for a coefficient lift of $\overline r$ of degree $e_0$ with unit leading coefficient.
If the residue characteristic divides $j$, then $r^j=t^j$ and $y$ has
degree at most $n-je_0\le n-2$, which was excluded.
Otherwise $j$ is a unit, and the identity
$j t^{i+j} f=(i+j)t^j y-i z$ gives $\deg f\le n-ie_0$.
Thus $i=e_0=1$, and its scalar degree still obeys
$\deg\overline f\le n-(i+j)e_0\le n-2$.

Write $r=uP_\gamma+e$ with $u\in A^\times$ constant and $e\in J[x]$.
At a surviving first term $\Lambda\mid\overline f$, hence
$f=P_\gamma g+\eta$ for some $g\in A[x]$, $\eta\in J[x]$.
Thus $ef=P_\gamma eg$. For $m=i+j\ge2$,
$r^mf=u^mP_\gamma^mf+mu^{m-1}P_\gamma^{m-1}ef$ is
divisible by $P_\gamma^2$. This excludes the progression.
Passing from cutoffs to a single set by scalar-degree gluing gives
\[
 1-\Delta_L(A)\le\frac{1-q^{-q}}{qQ}+\frac{q+1}{Q^2}.
\]

\paragraph{A matching of triples.}
Start with all $f$ in the rectangle with $\Lambda\nmid a=\overline f$.
Choose the first $\beta\in k$ in a fixed order with $a(\beta)\ne0$.
Put $P=P_\beta$, $s=n-\deg a\ge2$ and
\[
 c=\frac{-f_{n-1}}{2\operatorname{lc}(a)}\in J,\qquad
 r=P+c(P^s+1),\qquad y=rf,\quad z=r^2f .
\]
The operations on $c$ use the $k$-action on $J$. These formulas give
$\deg y,\deg z\le n$ and $y_n=-c\operatorname{lc}(a)$: the coefficient
of degree $n+1$ in $z=P^2f+2cP(P^s+1)f$ is
$f_{n-1}+2c\operatorname{lc}(a)=0$.
Discard $c=0$. Each remaining first term has degree $n-1$ and $y$ has
degree $n$. Require $z_n\ne0$ as well. If $s=2$ the reduction
$\overline z=(x-\beta)^2 a$ already has degree $n$. For $s>2$,
conditional on $a$ and $c$, the coefficient $f_{n-2}$ is uniform on
$J$ and occurs in $z_n$ with coefficient one; its occurrences in terms
multiplied by $c$ depend only on its reduction. Thus these two exclusions
cost at most $2/M$ of the starts.

At $\widetilde\beta$ one has $y(\widetilde\beta)=c a(\beta)\ne0$ and
$z(\widetilde\beta)=0$. If $\gamma\ne\beta$ is a root of $a$,
then $f(\widetilde\gamma)\in J$ and
\[
 y(\widetilde\gamma)=P(\widetilde\gamma)f(\widetilde\gamma),\qquad
 z(\widetilde\gamma)=P(\widetilde\gamma)^2f(\widetilde\gamma).
\]
The factor $P(\widetilde\gamma)$ is a unit. Conditional on $a$ and $c$,
$f(\widetilde\gamma)$ is uniform on $J$: its constant coefficient has
a free uniform $J$-part. Require all these evaluations to be nonzero,
at cost at most $q/M$. At other residue points $y,z$ have nonzero
reductions. Consequently $y$ has no root among the chosen representatives,
while $z$ has exactly one, $\widetilde\beta$.

One further condition identifies the start from $y$. For every simple root
$\gamma$ of $Y=\overline y$ test the equality in $J$
\[
 y(\widetilde\gamma)=
       -y_n\,\frac{(Y/(x-\gamma))(\gamma)}{\operatorname{lc}(Y)} .
\]
It holds at $\beta$. At any other simple root of $Y$, the right side is
fixed after conditioning on $a,c$, and the left side is uniform on
$J$, as just proved. Discard starts with a second successful test,
at cost at most $q/M$. All these costs use the union bound on the
original starts and require no independence between the tests.

We obtain disjoint triples. Degree distinguishes starts from both
other positions, and evaluation at the representatives distinguishes
middles from ends. A middle identifies $\beta$ by its unique
successful test. Then $a=\overline y/(x-\beta)$ and
$c=-y_n/\operatorname{lc}(a)$ recover $r$.
An end identifies $\beta$ by its unique root and gives
$a=\overline z/(x-\beta)^2$. Dividing in $A[x]$ by the monic $P$
yields $(z/P)(\widetilde\beta)=2c a(\beta)$ and hence $c$.
Finally $r$ is a non-zero-divisor, since
$r(P-c(P^s+1))=P^2$ is monic. Each inversion thus recovers at most one
start. The ordinary triple belongs to the allowed family for each $L\ge2$. Removing one member of each disjoint triple forces
\[
 1-\Delta_L(A)\ge
 \frac{1-q^{-q}}{qQ}
 \left(1-\frac{2+2q}{M}\right).
\]
Together the two bounds give the limit.

\clearpage
\section{A linear bound for stabilization in rank-two graph towers}
\label{sec:rank-two-onset}
Let $p$ be prime. A rank-two integral voltage tower attaches a vector of
$\mathbb Z^2$ to each oriented base edge, with negation on reversal. At
level $n$ the deck group is $(\mathbb Z/p^n\mathbb Z)^2$ and following
an edge adds its voltage to the deck coordinate. Write $v_n$ for the
$p$-adic valuation of the number of spanning trees. DuBose and
Vallières proved eventual polynomiality of $v_n$ in $p^n,n$
(\href{https://doi.org/10.5802/alco.304}{Theorem~6.2}).
The preceding \href{https://github.com/VladimirReshetnikov/ProveIt/blob/bd8864aa31430944f587f544213d562665f1d817/docs/incoming/Effective_Graph_Towers.zip}{ProveIt graph-tower report}
gives an algorithm in every rank and asks, in its first further research
question, for the growth of the worst possible onset as a function of
voltage bit length. We obtain a sharp order in rank two at $p=3$ for a
bouquet of five loops. The upper estimate works for all primes and
rank-two Laurent determinants. It uses the cyclotomic identity
$v_p(\zeta_{p^h}-1)=1/\varphi(p^h)$, also used in
\href{https://github.com/lsernau/grh-lean/blob/72fead102f663d5c7a407964a13dc0cf9d977207/grh/src/New/Congruences/Selmer/RootsOfUnity.lean}{the repository's specialization argument}.
The new ingredient is separation of the finitely many rational slopes
where pairs of Laurent monomials coincide.

\paragraph{Quantitative shell theorem.}
Let $F\ne0$ be an integral Laurent polynomial in two variables,
supported in $[-H,H]^2$, with $H\ge1$. Assume $F$ is nonzero at all
nontrivial pairs of $p$-power roots of unity. For $h\ge1$ write
$S_h$ for the sum of $v_p(F(x,y))$ over pairs of joint order exactly
$p^h$. Set
\[
 K=\lceil\log_p(8H^2+1)\rceil,\quad
 M=\lceil\log_p(2H+1)\rceil,\quad h_0=K+M+2.
\]
There exist rational $s_2,s_1,t_1,s_0$ such that for every $h\ge h_0$,
\[
 S_h=s_2p^{2h}+(s_1+t_1 h)p^h+s_0.
\]
All constants are independent of the level. The bound $h_0$ only
depends on the exponents and $p$.

\paragraph{Proof.}
Normalize $v_p(p)=1$ and write $e_h=(p-1)p^{h-1}$.
Dividing $F$ by the common $p$-power $p^g$ of its coefficients adds
$g(p^{2h}-p^{2h-2})$ back to the shell. Collect equal exponents and
assume henceforth that $F=\sum_i c_i X^{a_i}Y^{b_i}$ has a unit
coefficient. Fix compatible primitive roots $\zeta_h$. Simultaneous
raising to units modulo $p^h$ preserves the valuation, with each
orbit of size $e_h$. Its contribution is therefore
$\operatorname{ord}_h(F)=e_h v_p(F)$. Each orbit is uniquely
represented in one of two charts, $(\zeta_h,\zeta_h^z)$ for any $z$
modulo $p^h$, or $(\zeta_h^z,\zeta_h)$ with $p\mid z$. In the
second chart interchange the two exponent coordinates below.

The projected exponents in the first chart are $a_i+b_i z$. Take as
centers the distinct $p$-integral rational solutions of
$a_i+b_i z=a_j+b_j z$ lying in the chart. Write each center as a
reduced fraction $A/B$, with $p\nmid B$, $|A|,|B|\le2H$. The
difference of two distinct centers has nonzero numerator of absolute
value at most $8H^2$. Hence their disks $z\equiv A/B\pmod{p^K}$
are disjoint. In the restricted chart each disk is entirely inside
or outside the chart since $K\ge1$.

Outside these disks the projected exponents are pairwise distinct
modulo $p^J$, where $J=K+M$. In fact a nonzero slope difference
$b_i-b_j$ has valuation $<M$, and the distance factor
$z-(a_j-a_i)/(b_i-b_j)$ has valuation $<K$; this covers centers
outside the chart too. A nonzero constant difference has valuation
$<M$. For each fixed $z\bmod p^J$ in this complement let
$n_i\in[0,p^J)$ be the distinct exponent residues. The polynomial
$\sum_i c_i(1+T)^{n_i}$ is nonzero modulo $p$ by the distinct
degrees and the unit coefficient. Its first nonzero coefficient
occurs at $j<p^J$. For any lift of $z$ and $h\ge h_0$, the orbit
contributes exactly $j$. Indeed, modulo $(p,T^{p^J})$ the binomial
expansion of each monomial depends only on its exponent modulo $p^J$.
One may choose nonnegative representatives modulo $p^h$ in forming
the actual polynomial. At $T=\zeta_h-1$ a coefficient divisible
by $p$ has order at least $e_h\ge p^J$, and all later terms have
larger order than $j$. Each class contributes $j p^{h-J}$ to $S_h$.

Fix a center $A/B$. Applying the Galois automorphism
$\zeta_h\mapsto\zeta_h^B$ expresses the same valuation as the one of
\[
 P(\zeta_h,W),\qquad
 W=\zeta_h^{Bz-A},\qquad
 P(X,Y)=\sum_i c_i X^{B a_i+A b_i}Y^{b_i},
\]
up to multiplication by a monomial of valuation zero. Shift each
exponent axis by its minimum so that $P$ is a polynomial of degrees
$D\le8H^2$ in $X$ and $E\le2H$ in $Y$. The exponent transformation
is injective, so $P$ still has a unit coefficient. As $z$ ranges over
the disk modulo $p^h$, $W$ ranges exactly once over the roots of
order dividing $p^{h-K}$, since $B$ is a unit. Write $r$ for the
conductor exponent of $W$, with $r=0$ for $W=1$, and put $R=M+1$.

For $0\le r\le R$, use $e_0=1$ and write
$P(1+T,W)=\sum_{j=0}^{D} d_j(W)T^j$.
This polynomial is nonzero: its vanishing would give a torsion zero
of $F$ at the pair represented by $z$. Its nonzero coefficients have
valuations in $e_r^{-1}\mathbb Z$. Let $\mu_W$ be their minimum
and let $j_W$ be the first index attaining it. We have
$e_h/e_r>D$ (for $r\ge1$ the quotient is $p^{h-r}$ and for $r=0$
use $e_h\ge p^{h-1}$). Thus terms with larger coefficient valuation
cannot minimize. The unique minimum has order $e_h\mu_W+j_W$.
The set of all $W$ with $r\le R$ is fixed for every $h\ge h_0$.
Its total is a constant multiple of $p^h$ plus a constant.

For $R<r\le h-K$, put $k=h-r$. Then
$\operatorname{ord}_h(W-1)=p^k>D$ and $e_r>E+1$.
Expand
$P(X,1+U)=\sum_{i=0}^E P_i(X)U^i$, and take the smallest $l$
such that $P_l(X)$ is nonzero modulo $p$. Such an $l$ exists.
Let $d\le D$ be the order of $P_l(1+T)$ at $T=0$ modulo $p$.
Its value at $X=\zeta_h$ has order exactly $d$. Terms with $i<l$
have order at least $e_h=e_rp^k>lp^k+d$. Terms with $i>l$ have
order at least $(l+1)p^k>lp^k+d$. The orbit contribution is exactly
$lp^k+d$ for every root of conductor $p^r$. The contribution from
all these roots is
\[
 \sum_{k=K}^{h-R-1}e_{h-k}(lp^k+d)
 =l(h-R-K)e_h+d(p^{h-K}-p^R).
\]
Adding the finitely many disks in each chart, their complements,
and the removed content proves the shell formula.

\paragraph{Consequence for graph towers.}
Let $F$ be the determinant of the Laurent voltage Laplacian of a
finite connected base graph, and assume the finite lifts are
connected. Diagonalizing by the characters of the deck group splits
the lifted Laplacian into blocks. The trivial character yields the
base Laplacian; connectedness gives invertible blocks at every
nontrivial character, since the lifted Laplacian has nullity one.
Applying the matrix-tree identity to the base and the lift gives
\[
 v_n=v_0-2n+\sum_{h=1}^{n} S_h .
\]
Summing the shell formula shows that from $n=h_0$ onward
\[
 v_n=A p^{2n}+(B+C n)p^n+D' n+E'
\]
for rational constants. For a bouquet the support of its Laurent
Laplacian consists of zero and the signed voltage vectors. Thus if
the absolute values of the coordinates are at most $2^b$, the
first stable level is at most $3b\log_p 2+O_p(1)$. For a fixed base
graph the coordinates in the determinant support are bounded by
the vertex count times the largest voltage coordinate, by determinant
expansion. The same bit-length estimate follows with a base-dependent
additive constant. Five consecutive values beyond this upper bound
suffice to recover the law: the evaluation matrix on
$1,n,p^n,np^n,p^{2n}$ is nonsingular since these sequences solve the
order-five recurrence with characteristic polynomial
$(x-1)^2(x-p)^2(x-p^2)$ and are linearly independent. Checking the
finitely many preceding levels then gives the least onset. This is a
bound on conductor precision; enumerating all the characters there
can still be expensive.

\paragraph{Matching lower bound at $p=3$.}
The preceding graph report gives five-loop bouquets with directions
$3^s e_1,e_2,e_1,e_1,e_1$ for $s\ge2$.
For completeness let $A(U)=2-U-U^{-1}$. At joint conductor $3^j$
the sum $A(U)+A(V)$ has valuation exactly $2/e_j$: its normalized
quadratic coefficient in $\zeta_j-1$ is $-(a^2+b^2)$ for exponents
$a,b$ not both divisible by $3$, which is nonzero modulo $3$.
Consequently $A(X^{3^s})+A(Y)$ dominates $3A(X)$ off the
kernel of $(X,Y)\mapsto(X^{3^s},Y)$, since its valuation is at most
one. Put $t=\min(n,s)$. The image is $C_{3^{n-t}}\times C_{3^n}$
and each fiber has size $3^t$. Its order-$3^j$ shells have sizes
$3^{2j}-3^{2j-2}$ for $j\le n-t$ and $3^{n-t}e_j$ for $n-t<j\le n$.
The kernel contributes $3^t-1+2t$, using
$\prod_{X^{3^t}=1,\ X\ne1}(1-X)=3^t$.
Adding these contributions and subtracting $2n$ gives
\[
 v_n=2t3^n+4(3^n-3^t)+(3^t-1)+2t-2n.
\]
At $n\ge s$ this is the eventual polynomial obtained by substituting
$t=s$. At $n=s-1$ its actual value exceeds the extrapolated polynomial
by $4\cdot3^{s-1}-2$. The least onset is therefore $s$. With
$3^s\le2^b$, the worst-case onset for this fixed prime and loop count
grows as $\Theta(b)$. The sharp leading constant and higher ranks
remain questions.

\paragraph{Sharp order at every prime for fixed loop count.}
The coefficient of $p^{2n}$ in the graph formula is exactly the
common $p$-adic content $g=\min_i v_p(c_i)$ of its Laurent
determinant, in the sense of the valuation of the content. Indeed,
restoring the removed power contributes exactly $g(p^{2n}-1)$
to the cumulative sum. Every other summand grows as $O(np^n)$.
This yields a second lower bound that uses no restriction on $p$.
On a bouquet of $2p+1$ loops put $p$ copies each of $e_1,e_2$ and
one step $p^s e_1$, with $s\ge1$. Compare with the same assignment
where the final step is zero. Both towers are connected, and the
graphs agree through level $s$. The zero-step determinant is
$p(A(X)+A(Y))$ and has content valuation one. The extra summand
$A(X^{p^s})$ changes the content valuation to zero because its
two nonconstant exponents are new. The two eventual laws have
different $p^{2n}$ coefficients. For $s$ large enough that
$s-4\ge h_0(1)$, the zero-step law is already exact throughout
$[s-4,s]$. If the other law were exact from $s-4$ onward the
two would agree on these five consecutive levels, which the
recurrence argument excludes. Its least onset is at least $s-3$.
Taking $s=\lfloor\log_p(2^b)\rfloor$ gives worst-case onset
$\Theta_p(b)$ on this one fixed bouquet for every prime. Padding
with zero-voltage loops gives the same order for any larger fixed
loop count. The five-loop ternary calculation above gives a
smaller loop count at that prime and identifies the exact onset.

\paragraph{Classification by number of loops.}
In fact the threshold loop count can be found exactly. Let $L_{p,m}(b)$
be the largest least onset over connected rank-two towers on one bouquet
of $m\ge2$ loops, with coordinates bounded by $2^b$ in absolute value.
Repeated directions and zero voltages are allowed. As $b\to\infty$,
\[
 \boxed{\quad L_{p,m}(b)=
 \begin{cases}
 O_p(1),&2\le m<\max(3,p),\\
 \Theta_p(b),&m\ge\max(3,p).
 \end{cases}\quad}
\]
Constants here can also be taken independent of $m$ in the upper
bound. For odd $p$ and $m<p$, the onset is at most one.

To see the bounded range for odd $p$, take a projective direction
modulo $p$. The projected step residues are not all zero because
connectedness forces the steps to span $\mathbb F_p^2$. In the group
ring $\mathbb F_p[Z]/(Z^p-1)$ take any nonzero residue $a$ occupied
by a step, and let $j$ count the steps in $\{a,-a\}$. The
coefficient of $Z^a$ in the projected Laplacian is $-j\ne0$
since $a\ne-a$ and $1\le j\le m<p$. Thus its expansion at $Z=1$
has first surviving order $<p$. Since $e_h\ge p$ for $h\ge2$,
every orbit of depth $h$ lifting that projective direction has exactly
that integral order. Every residue direction has $p^{h-1}$ lifts,
so $S_h$ is a fixed multiple of $p^{h-1}$ for $h\ge2$. Summation
from level one proves the claim.
For $p=2,m=2$ the two steps form an invertible basis modulo $2$,
hence modulo $2^n$ for every $n$. A change of coordinates gives the
same square grid at every fixed level for all choices of voltages,
so the support bound $h_0$ evaluated at $H=1$ suffices.

For the lower range it is enough to work at the threshold value
and pad with zero voltages. We use a test supplied by the shell proof.
If a Laurent polynomial is primitive modulo $p$, the coefficient of
$np^n$ in the cumulative formula is positive precisely when some
nonzero rational direction $(u,v)$ satisfies
\[
                F(Z^u,Z^v)\equiv0\pmod p
                \quad\text{as a Laurent polynomial in }Z,
\]
where $(u,v)$ is taken to be a primitive integer representative.
Indeed the coefficient of $h p^h$ in $S_h$ is the sum of the
nonnegative disk contributions $l(p-1)/p$. A disk contributes
positively exactly when $P(X,1)$ vanishes modulo $p$. This is the
displayed identity after the monomial shift and the change of
coordinates defining the disk. Conversely any such identity must
group at least two distinct exponents with the same projection
because there is a unit coefficient. This yields one of the
rational centers, using the chart where the direction has a unit
pivot. Summing $hp^h$ gives a positive multiple of $np^n$
plus a multiple of $p^n$ and a constant. This proves the test.

Suppose $m=p$ is odd. Use $p-2$ copies of $(1,0)$, one copy
of $(1,1)$ and one copy of $(1+\delta,2)$, first with $\delta=0$
and then with $\delta=p^s>4$. In both cases the content is a unit
and the first two directions ensure connectedness. In the first
case projection along $(u,v)=(1,0)$ makes all $p$ exponents equal
to one, so the test is positive. For the second case a vanishing
projection must have all $p$ absolute projected exponents equal
to one another and nonzero. Indeed an occupied positive absolute
frequency $d$ contributes the negative of its multiplicity to the
coefficient of $Z^d$. This multiplicity must equal $p$, and a
nonzero direction cannot annihilate every step. Thus $u\ne0$ and
$|u+v|=|u|$, so $v=0$ or $v=-2u$. The last projected absolute
value is now $|(1+\delta)u|$ or $|(\delta-3)u|$, both different
from $|u|$. The $np^n$ term vanishes after perturbation.

For $p=2,m=3$ use the three steps
$(1,0),(1,1),(\delta,1+2\delta)$ with $\delta=0$ or $2^s$.
Again the content is a unit. In characteristic two every positive
absolute projected frequency must occur an even number of times.
Three steps spanning the plane therefore give a zero Laurent
polynomial exactly when one projected value is zero and the other
two have equal nonzero absolute value. A direction annihilating
one step projects the other two in proportion to their determinants
with that step. For $\delta=0$ the three pairwise absolute
determinants all equal one, and such a direction exists. For $\delta>0$
the determinants are $1,1+\delta,1+2\delta$, all distinct.
Thus this pair also has different $np^n$ coefficients.

In both constructions the perturbed tower agrees exactly with the
fixed initial tower up through level $s$. Once the fixed initial
tower has stabilized by $s-4$, the five-consecutive-values argument
forces onset at least $s-3$ for the perturbed tower. The heights
are at most $1+2p^s$, allowing $s=b\log_p 2+O_p(1)$. Together with
the upper bound this proves the classification. It leaves open the
optimal leading constant within the linear-growth range.

\clearpage
\section{A bipartite family with all zeros on the canonical line}\label{sec:ehrhart}
For a graph $G$ on $v$ vertices let $P_G$ be the convex hull of
$\pm(e_i-e_j)$ over its edges, in the lattice $\mathbb Z^v\cap\{\sum x_i=0\}$.
Write $E_{m,n}(k)=|kP_{K_{m,n}}\cap\mathbb Z^{m+n}|$ for its Ehrhart polynomial.
The degree is $d=m+n-1$. Higashitani, Kummer and Michałek asked whether all zeros lie
on $\Re z=-1/2$ and whether decreasing a largest part gives interlacing, for
complete multipartite graphs (Conjecture~4.10, \href{https://arxiv.org/abs/1612.07538}{arXiv:1612.07538}).
A recent public manuscript reports a counterexample with six parts of size two
(\href{https://itsdavis.com/solved-conjectures/results/hkm-ehrhart/}{Cinna Davis, proof and verification}).
It leaves the complete bipartite question intact. Kölbl gave a conditional
result for $K_{4,n}$ (\href{https://arxiv.org/abs/2404.02136}{arXiv:2404.02136}, Theorem~3.3).
Carptopus then proved the line property through smaller part five using positive Jacobi
matrices (\href{https://doi.org/10.5281/zenodo.22486582}{2026 preprint}).
These are direct antecedents of the construction below.

\paragraph{Finite-width theorem.}
For $2\le m\le 32$ and every integer $n\ge m$, all the zeros of $E_{m,n}$
are simple and lie on $\Re z=-1/2$. The ordinates of the zeros strictly alternate with
those for $E_{m,n-1}$. Width $32$ is the last width for which we certify a positive Bezoutian.
The proof uses no zero-location theorem for analytic $L$-functions.

\subsection*{Reduction to two short polynomials}
Put $r=m-1$, $d=n+r$, and let $C_k$ denote the Ehrhart polynomial of a $k$-dimensional
cross-polytope. Its Ehrhart series is $(1+t)^k/(1-t)^{k+1}$.
Set $Q_k(y)=k!\,i^{-k}C_k((iy-1)/2)$. The cross-polytope recurrence
gives $Q_0=1$, $Q_1=y$, $Q_{k+1}=yQ_k-k^2Q_{k-1}$.
In particular, $Q_k$ is the characteristic polynomial of a real symmetric
tridiagonal matrix with consecutive off-diagonal entries $1,2,\ldots,k-1$.
The formula of Higashitani, Jochemko and Michałek for the $h^*$-polynomial
(\href{https://arxiv.org/abs/1807.07678}{arXiv:1807.07678}, Theorem~4.1)
has coefficients $\gamma_j=\binom{2j}{j}\binom{r}{j}\binom{n-1}{j}$ in the
basis $t^j(1+t)^{d-2j}$. Use $4t=(1+t)^2-(1-t)^2$ in the Ehrhart series.
Binomial expansion expresses $4^rE_{m,n}$ as
$\sum_{h=0}^r(-1)^h c_h C_{d-2h}$, where
$c_h=\sum_{j=h}^r4^{r-j}\gamma_j\binom{j}{h}$.
Thus the rotated polynomial has the exact expression
\[
 R_n(y)=4^r d!\,i^{-d} E_{m,n}((iy-1)/2)
       =\sum_{h=0}^r c_h(d)_{2h}Q_{d-2h}(y)
       =A_n(y)Q_n(y)-B_n(y)Q_{n-1}(y),
\]
where $(d)_k=d(d-1)\cdots(d-k+1)$. Expand each $Q_{d-2h}$ from
$Q_n,Q_{n-1}$ by forward or backward recurrence to define $A_n,B_n$.
Their $y$-degrees are at most $r,r-1$. $A_n$ has degree $r$ with
positive leading coefficient $c_0$. Their parities are $r,r-1$ respectively.
The construction also works at $n=m-1$, including a possible $Q_0$.

\subsection*{The positive form and the interlacing determinant}
Let $A=A_n$, $B=B_n$. Form the real symmetric matrix $H$ of coefficients in
$u^iv^j$ ($0\le i,j<r$) of the Bezoutian
\[
 \frac{A(u)B(v)-B(u)A(v)}{u-v}.
\]
Suppose $H$ is positive definite. Substitution $v=\bar u$ shows that
$A,B$ have no common zero and that $\Im(B(u)/A(u))<0$ for $\Im u>0$
where the quotient is defined. In fact $A(u)$ cannot vanish there, since
the diagonal Hermitian form at $(1,u,\ldots,u^{r-1})$ would vanish.
On the upper half-plane $Q_n/Q_{n-1}$ has positive imaginary part by
induction using $Q_{k+1}/Q_k=y-k^2/(Q_k/Q_{k-1})$.
The two ratios cannot coincide, so $R_n$ has only real zeros.
On the real axis the diagonal limit gives $W=A'B-AB'>0$.
Also $W_n=Q_n'Q_{n-1}-Q_nQ_{n-1}'>0$, since $W_1=1$ and
$W_n=Q_{n-1}^2+(n-1)^2W_{n-1}$ for $n\ge2$.

Set $C=y A_{n+1}-B_{n+1}$ and $D=n^2 A_{n+1}$, so that
$R_{n+1}=CQ_n-DQ_{n-1}$. The polynomial
\[
 T_n(y)=AD-BC=n^2 A_n A_{n+1}-y B_n A_{n+1}+B_n B_{n+1}
\]
is even in $y$. At every real root $t$ of $R_n$ there is a nonzero
$\lambda$ with $(Q_n,Q_{n-1})(t)=\lambda(B,A)(t)$.
Both vectors are nonzero by the recurrence and positive definiteness.
Differentiation gives $R_n'(t)=(\lambda^2W(t)+W_n(t))/\lambda$,
which proves simplicity. Moreover
\[
 \frac{R_{n+1}(t)}{R_n'(t)}
 =-\frac{\lambda^2 T_n(t)}{\lambda^2 W(t)+W_n(t)} .
\]
When $T_n$ is positive on the real axis, signs force a root of
$R_{n+1}$ in each interior gap and one in each exterior gap, since
both leading coefficients are positive. The degree difference is one,
so this accounts for all its roots, with multiplicity one.

\subsection*{Exact certification for unbounded $n$}
Only short polynomials involving $r$ occur in the two positivity tests.
The following details explain how finite calculations prove the universal
quantifier over $n$. In performing the backward recurrence, multiply
both coordinates by $S(n)=\prod_{j=1}^{r-1}(n-j)^2$.
Backward steps divide by distinct squared linear factors of $S$.
In the forward direction at most $r$ steps increase the degree
in $n$ by at most $2r$. In the backward direction the coordinates
have denominators dividing $S$ and, after multiplication by $S$,
degrees at most $2r$ (each step is $[yq_{k+1}-q_{k+2}]/(n+k+1)^2$,
where indices are relative to $n$). Thus each cleared coordinate in the
mixture has degree at most $7r$ in $n$:
$2r$ allows for $S$, $2r$ for forward recurrence, $r$ for $c_h$,
and $2r$ for $(d)_{2h}$. At $7r+1$ consecutive integers beginning at
$m$, the checker verifies that the coordinates divided by $S$ agree with
polynomials of degrees at most $r$ for $A_n$ and $r+1$ for $B_n$, by taking
forward differences. Multiplying these interpolants by $S$ and applying the
polynomial identity principle proves the degree bounds identically.

The Bezoutian splits into even and odd blocks by parity. Each entry has
degree at most $2r+1$ in $n$. For the leading minor of order $j$ of
either block, the checker computes its exact integer values at
$n=m,m+1,\ldots,m+(2r+1)j$ by fraction-free elimination.
All forward differences at $m$, through its actual degree, are positive.
For every integer $s\ge0$, Newton's identity
$F(m+s)=\sum_h\binom{s}{h}\Delta^h F(m)$ now proves strict positivity.
Sylvester's criterion applies to both blocks. At $n=m-1$ the two
blocks are also tested directly.

Finally, write $T_n(y)=U(n,y^2)$. Its degree in $n$ is at most $2r+2$.
The checker obtains $\Delta_n^h U(m,z)$ from $2r+3$ rows. For each nonzero
polynomial in $z$ it verifies a positive constant term and equal numbers
of Sturm sign changes at $0$ and $+\infty$. Hence each is positive
on $[0,\infty)$, and so is $U(m+s,z)$ for integral $s\ge0$ by
Newton's identity (the zeroth row is positive). The same Sturm test at
$n=m-1$ completes the square boundary case.
The scripts \texttt{verify\_ehrhart.py} and \texttt{verify\_interlacing.py}
perform the integer calculations. Each accepts the first and last widths
as arguments. Run both with \texttt{2 32}. Width one follows at once from
the cross-polytope recurrence. A proof for all complete bipartite graphs
would require a uniform argument controlling the forms or another way to
localize the roots as $r$ grows.

\subsection*{Interlacing when the smaller part increases}
For $2\le m\le32$ and $n\ge m$, the ordinates of the zeros of
$E_{m,n}$ strictly interlace those of $E_{m+1,n}$. Thus the line
property holds also for every $K_{33,n}$ with $n\ge33$. Consecutive-$n$
interlacing within width $33$ remains outside this certificate.
In particular, the interlacing of $E_{3,n}$ and $E_{4,n}$ follows
unconditionally for $n\ge3$. The cases $n=1,2$ of this particular
interlacing follow from the cross-polytope recurrence and the
larger-part interlacing at width two, respectively.

To see this, keep the pair $Q_n,Q_{n-1}$ fixed, and use subscripts
$(m,n)$ for the short coordinates. In the residue formula take
$C=A_{m+1,n}$ and $D=B_{m+1,n}$. The determinant
\[
 V(n,y^2)=A_{m,n}(y) B_{m+1,n}(y)-B_{m,n}(y) A_{m+1,n}(y)
\]
has even parity. Its degree in $n$ is at most $2r+2$, as the termwise
degree bounds for its factors are $(r,r+2)$ and $(r+1,r+1)$.
The script \texttt{verify\_smaller\_part.py} tests the coordinate bounds
for width $m+1$ by the same interpolation identity; its upper endpoint
includes width $33$. It certifies $V(m+s,z)>0$ for every integer
$s\ge0$ and real $z\ge0$ by applying the Sturm test to the Newton
coefficient polynomials in $s$. The residue sign argument gives the
claimed interlacing. This argument needs the positive Bezoutian only
in the lower width. Run the script with \texttt{2 32}.

\subsection*{A worked example and an independent check from the polytope}
For $K_{6,6}$, dividing both short coordinates by their common factor $132$ gives
\[
 \begin{aligned}
 A(y)&=693y^5-29984y^3+175507y,\\
 B(y)&=19656y^4-442620y^2+881100 .
 \end{aligned}
\]
Thus positivity of a five-by-five real form suffices for the degree-eleven Ehrhart polynomial in this example. Its first four lattice-point counts, at dilations $0,1,2,3$, are $1,73,1915,26547$.

There is a direct way to check these counts. Split a lattice vector into its coordinates $(u,v)$ on the two sides of the graph. Its membership in $kP_{K_{m,n}}$ is equivalent to
\[
       {\textstyle\sum} u_i+{\textstyle\sum} v_j=0,\qquad
       \max(\|u\|_1,\|v\|_1)\le k.
\]
Indeed, a weighted sum of signed edges is a matrix with row sums $u$, column sums $-v$, and cost equal to the sum of the absolute entries. Both marginal norms bound that cost from below. For the converse let $T$ be the larger marginal norm and let $s=\sum u_i$. On each side choose nonnegative positive and negative margins with total positive mass $(T+s)/2$ and total negative mass $(T-s)/2$. Starting from the ordinary positive and negative parts, add the same padding to both where needed. Nonnegative matrices with each pair of prescribed margins exist by successively matching supplies and demands. The difference of the two matrices realizes the desired vector with cost at most $T$. Real edge weights suffice for membership in the polytope.

The independent script \texttt{check\_instances.py} counts integer vectors by their sum and $\ell^1$ norm on each side. It recovers the entire polynomial from the required $d+1$ integer values for $K_{2,2},K_{3,5},K_{6,6},K_{7,9}$, a total of 40 comparisons. The proof for unbounded part sizes uses the Newton and Sturm certificates above.

\clearpage
\section{When the Lubin--Tate branch bound is sharp}
\label{sec:branch-equality}

This calculation answers the question about cancellation left by the
\href{https://github.com/VladimirReshetnikov/ProveIt/blob/5cc267af75a635199ac5d93a7d49cbce521daefc/docs/incoming/grh_transfers_cycle_signings_and_lubin_tate.zip}{earlier branch-interpolation report}.
Its source in the Lean study was the use of an extremizer to establish the
sharpness of a derivative bound
(\href{https://github.com/lsernau/grh-lean/blob/72fead102f663d5c7a407964a13dc0cf9d977207/grh/src/New/Derivatives/Sharpness/OneCurve.lean\#L165-L183}{the one-curve lemma}).
The application concerns finite precision for the branch-evaluation method
in Berger and Sprang's \href{https://arxiv.org/abs/2502.18053}{Lemma~4.2 and Proposition~4.4}.

Let $\mathcal O$ be the ring of integers of a finite extension of
$\mathbb Q_p$, with uniformizer $\pi$, valuation $v(\pi)=1$ and residue
field of size $q$. Fix a Lubin--Tate formal group for $\pi$.
Its endomorphism $[a](X)$ has linear term $aX$ and
$[\pi](X)\equiv X^q\pmod\pi$ in a Lubin--Tate coordinate.
Let $A=\mathcal O/\pi^r$, $r\ge1$.
For a finite set $B\subset\mathcal O$ of at least two distinct directions
and $b\in B$, put
\[
 h_c=v(b-c),\quad h=\max_{c\ne b}h_c,\qquad
 \omega=\sum_{c\ne b}q^{h_c},\qquad
 \Delta=\prod_{c\ne b}([b](X)-[c](X)).
\]
Write $t$ for the least integer such that $X^t/\Delta\in A[[X]]$.
The previous report bounds $t\le\omega+(r-1)(q^h-1)$.
The branch cutoff is $1+t$ when recovering all coefficients of total degree
less than $|B|$. Indeed the Lagrange numerator
$P_b=\prod_{c\ne b}(Y-[c](X))$ has total order $|B|-1$ and leading
$Y$-coefficient one. If all observations vanish at their indicated
precisions, each term $H(X,[b](X))P_b/\Delta$ of the interpolation
remainder has total order at least $|B|$. The discarded multiple of
$\prod_c(Y-[c](X))$ has the same property. Division by this monic
distinguished polynomial is valid over the complete local ring $A[[X]]$.
Conversely $(X^t/\Delta)P_b$ has a nonzero constant-in-$X$ coefficient of $Y^{|B|-1}$
and vanishes on every other branch. Its value on branch $b$ is $X^t$.

\paragraph{Equality criterion.}
The bound is an equality if $r=1$ or $h=0$. For $r>1$ and $h>1$ it is
always strict. For $h=1$, set $N=\#\{c\ne b:h_c=1\}$ and define
\[
 j_* =\max\left\{0\le j<r:
         j+v\!\binom{N+j-1}{j}<r\right\}.
\]
Here $v$ acts on the indicated positive integer, with the normalization
above. The exact pole order is
\[
                    t=\omega+(q-1)j_* .
\]
Consequently at height one equality in the earlier bound holds exactly
when $p$ does not divide $\binom{N+r-2}{r-1}$.

\paragraph{Proof.}
The difference $[b]-[c]$ is $[\pi^{h_c}]$ times a unit in $\mathcal O[[X]]$.
In fact $F(U,V)-U$ equals $V$ times a unit, where $F$ is the formal
addition law. Substitute $U=[c]$ and $V=[b-c]$, and use
$[b-c]=[u]\circ[\pi^{h_c}]$ with $u$ a unit.
Multiplication by a power-series unit preserves the least exponent with
nonzero coefficient in $A((X))$, including when that coefficient is a
zero divisor. Hence all differences at the same height can be replaced
by identical factors.

Height zero gives $[1](X)=X$. For $a\ge1$, put $s=q^a$.
Since $[\pi^a](X)\equiv X^s\pmod\pi$ and the linear term is
$\pi^aX$, write
\[
            [\pi^a](X)=X^s U_a(X)+\pi X V_a(X),\qquad
            U_a(0)\in\mathcal O^\times,\quad V_a(0)=\pi^{a-1}.
\]
Take the terms of degrees below $s$ into $\pi XV_a$.
Inverting the factor in $A((X))$ amounts to expanding
$(1+\pi X^{1-s} V_a/U_a)^{-1}$ through power $r-1$ and multiplying
by $X^{-s}U_a^{-1}$. In the product of these expansions at most
$r-1$ powers of $\pi$ can contribute. Each can add at most $q^h-1$
to the pole order. This gives the earlier bound. Modulo $\pi$ the
reciprocal has pole order $\omega$, which proves equality if $r=1$.
If $h=0$, all factors are $X$ times units at every precision.

Suppose $h>1$ and $r>1$. A contribution at the extreme exponent would
require $r-1$ factors of height $h$, counted with repetition, and would
use only their constant terms $V_h(0)/U_h(0)$. Each such constant is
divisible by $\pi$. The coefficient therefore vanishes modulo $\pi^r$.

Finally suppose $h=1$. Up to a power-series unit the reciprocal is
\[
 X^{-\omega}
 \sum_{j=0}^{r-1}(-1)^j\binom{N+j-1}{j}
       \pi^j X^{-j(q-1)}(V_1/U_1)^j .
\]
The power series $V_1/U_1$ is a unit. Each summand whose scalar survives has
least exponent $-\omega-j(q-1)$. The largest surviving $j$ thus
determines the least exponent in the whole sum. This proves the formula.

\paragraph{Exact answer modulo $\pi^2$ at arbitrary contact height.}
For each $a\ge1$ put $N_a=\#\{c\ne b:h_c=a\}$.
If all $N_a$ are divisible by $p$, then $t=\omega$ at $r=2$.
Otherwise let $a_*$ be the largest $a\ge1$ for which $p\nmid N_a$. Then
\[
                  t=\omega+(q-1)q^{a_*-1}\qquad(r=2).
\]
To see this, write $[\pi](X)=X^q+\pi W(X)$, where $W(0)=0$ and
$W'(0)=1$. Reduction of $[\pi^a]=[\pi]\circ[\pi^{a-1}]$ gives
$[\pi^a](X)\equiv X^{q^a}+\pi W(X^{q^{a-1}})\pmod{\pi^2}$:
adding a multiple of $\pi$ to the argument leaves its $q$th power
unchanged modulo $\pi^2$, since $p\mid q$. Up to a power-series unit,
the reciprocal of $\Delta$ is therefore
\[
 X^{-\omega}\left(1-
 \pi\sum_{a\ge1}N_a X^{-q^a}W(X^{q^{a-1}})\right)\pmod{\pi^2}.
\]
Each surviving term indexed by $a$ has its first power at
$-(q-1)q^{a-1}$ inside the parentheses. These exponents are distinct,
and the largest surviving height determines $t$.

For the multiplicative formal group over $\mathbb Z_2$ take $r=2$.
At $b=0$ for $B=\{0,2,6\}$ one has $N=2$, $\omega=4$ and $j_*=0$,
so the required cutoff is $5$; the former bound gave $6$.
For $B=\{0,4\}$ at $b=0$, the polynomial
$[4]=(1+X)^4-1$ reduces to $X^4+2X^2$ modulo $4$.
Its reciprocal is $X^{-4}-2X^{-6}$, giving $t=6$ against the bound $7$.


\clearpage
\section{Further research questions}
\begin{enumerate}
\item For finite local rings with $\mathfrak m^3=0$ and $\mathfrak m^2\ne0$, which invariant controls the loss of density? The elimination in \eqref{eq:eliminate} then leaves a term in the second radical layer. Determining whether its rank or its multiplication law enters the criterion is a first case beyond Theorem~\ref{thm:local}.
\item For a fixed prime and fixed loop count in the linear-growth range, what is the sharp coefficient of the voltage bit length? In rank three the exceptional directions can meet along a positive-dimensional set. An effective bound there would require controlling such intersections.
\item For complete bipartite edge polytopes, are the two parity blocks of the Bezoutian positive in every width? A representation of their minors as sums with positive coefficients could replace the width-by-width certificates and give a route to interlacing for unrestricted part sizes.
\item Can one give a closed expression for the pole order at every precision and arbitrary mixed contact heights? The calculation modulo $\pi^2$ detects the largest multiplicity surviving in the residue field. Higher precision also detects interactions between heights.
\end{enumerate}

\begin{thebibliography}{9}\small
\bibitem{lean} L. Sernau and collaborators, \emph{GRH for Dirichlet $L$-functions, in Lean 4}, based on the proof of L. Alpöge with Claude. \href{https://github.com/lsernau/grh-lean/tree/72fead102f663d5c7a407964a13dc0cf9d977207/grh}{Pinned readable edition}, \href{https://github.com/lsernau/grh-lean/blob/72fead102f663d5c7a407964a13dc0cf9d977207/grh/GUIDE.md}{reading guide}.
\bibitem{asada} M. Asada, E. Fourakis, S. Manski, N. McNew, S. J. Miller and G. Moreland, \emph{Subsets of $\mathbb F_q[x]$ free of 3-term geometric progressions}, Finite Fields Appl. \textbf{44} (2017), 135--147. \href{https://doi.org/10.1016/j.ffa.2016.10.002}{Published paper}, \href{https://arxiv.org/pdf/1512.01932}{questions in \S6 of the preprint}.
\bibitem{dhj} D. H. J. Polymath, \emph{A new proof of the density Hales--Jewett theorem}, Ann. of Math. \textbf{175} (2012), 1283--1327. \href{https://doi.org/10.4007/annals.2012.175.3.6}{Article}.
\bibitem{products} ProveIt, \emph{Cyclotomic valuations in graph towers and progression-free sets over finite rings}, \prev{1667e16cee9d4a68961e201e7288b8ecdac895c6}{grh_method_transfers_graph_towers_and_finite_rings.zip}{source packet}, especially \texttt{rings.tex} and further question 4. See also the \prev{5ed40f562710fcc2770bbb463dba528a6727f7ae}{grh_method_transfers_research.zip}{earlier local-ring bounds}.
\bibitem{graphprior} ProveIt, \emph{When does the spanning-tree formula start?} \prev{bd8864aa31430944f587f544213d562665f1d817}{Effective_Graph_Towers.zip}{Source packet}, first further research question.
\bibitem{branchprior} ProveIt, \emph{Cycle signings and Lubin--Tate branch interpolation}. \prev{5cc267af75a635199ac5d93a7d49cbce521daefc}{grh_transfers_cycle_signings_and_lubin_tate.zip}{Source packet}, branch-precision section.
\bibitem{fields} ProveIt, \href{https://github.com/VladimirReshetnikov/ProveIt/tree/main/docs/incoming}{incoming reports}. For the field results compare \texttt{grh-methods-polynomial-density.zip} and \texttt{grh-progression-densities-and-cubic-prime-race.zip}. Further primary references accompany the graph, Ehrhart and formal-group arguments.
\end{thebibliography}
\end{document}
